\documentclass[dvipdfmx]{jsarticle}

\usepackage{amsmath,amssymb}
\usepackage{bm}
\usepackage{graphicx}
\usepackage{here}
\usepackage{url}
\title{画像処理実験レポート}
\author{}
\date{}

\begin{document}

\maketitle


\section{目的}

デジタル画像に対する処理であるフーリエ変換と離散コサイン変換、畳み込みによる低域通過

フィルタリングと高域通過フィルタリングを実験により確認し、その処理方法とデジタル画像にお

ける周波数特性について理解する。

\section{原理}

\subsection{畳み込み}

古典的な画像処理手法は、所望の周波数特性を有するフィルタを用いたフィルタリングで実現さ

れている。フィルタは、二次元信号であり各要素はフィルタ係数を表している。画像 $X$ に対する

フィルタ $F$ によるフィルタリングは畳み込み演算 $X \ast F$ で実現され、以下の式で定義されている。



\begin{equation}
y_{ij} =
\sum_{m=-k_i}^{k_i}
\sum_{n=-k_j}^{k_j}
x_{(i+m)(j+n)} \times f_{-m,-n}
\tag{1}
\end{equation}

ここで、$y_{ij}$ と $x_{ij}$、$f_{ij}$ は出力画像 $Y$ と $X$、$F$ の $(i,j)$ 要素を表し、$2k_i+1$

と $2k_j+1$ は $F$ の縦横方向の長さをそれぞれ表している。つまり $F$ は、簡単のために縦横奇数の

長さを持つと仮定しており、中心座標を $(0,0)$ として中心から上下左右にそれぞれ $k_i$ と $k_j$

の要素数を持つことを表している。

\subsection{周波数解析}

離散時間信号の周波数解析には離散フーリエ変換（DFT: Discrete Fourier transform）を用い

るのが一般的である。長さ $N$ を持つ一次元信号 $x$ の DFT は以下の式で定義されている。

\begin{equation}
X(\omega)=
\sum_{n=0}^{N-1}
x_n \exp(-j\omega n)
\tag{2}
\end{equation}

ここで、$x_n$ は $x$ の $n$ 番目要素を表し、$\omega = 2\pi k/N$ である。定義より、$X(\omega)$ は

$2\pi$ の周期関数となり、$x$ の要素が実数値である場合、その周波数振幅特性 $|X(\omega)|$ と位相特性

$\angle X(\omega)$ は偶関数と奇関数になる。よって、$|X(\omega)|$ は $[0,\pi]$ の範囲を確認すれば良い。

二次元信号である画像 $X$ の周波数解析には二次元 DFT を用いる。二次元 DFT は、式(2)をそ

のまま二次元に拡張した以下の式で定義されている。

\begin{equation}
X(\omega_i,\omega_j)=
\sum_{m=0}^{M-1}
\sum_{n=0}^{N-1}
x_{mn}
\exp\{-(j\omega_i m+j\omega_j n)\}
\tag{3}
\end{equation}

ここで、各変数は式(1)と(2)と同様に定義されている。$|X(\omega_i,\omega_j)|$ は、$|X(\omega)|$

と同様に $2\pi$ の周期関数であり原点対称性を持つが、四象限対称ではないので、各変数の

$[-\pi,\pi]$ の範囲を確認する必要がある。

\subsection{アップサンプリングとダウンサンプリングと周波数振幅特性}

長さ $N$ を持つ一次元信号 $x$ に対して、各要素間に 0 を挿入して信号長を大きくする処理をアッ

プサンプリングと呼び、等間隔に要素を間引いて信号長を小さくする処理をダウンサンプリングと

呼ぶ。$x$ を倍率 $\gamma$ でアップサンプリングした結果を $y$ とすると以下の式で定義されている。

\begin{equation}
y_n =
\begin{cases}
x_m & \mathrm{if}\ \mathrm{mod}(n,\gamma)=0 \\
0 & \mathrm{otherwise}
\end{cases}
\tag{4}
\end{equation}

ここで、$\gamma$ は一般的に 1 より大きい整数であり、$m=\mathrm{div}(n,\gamma)$、$n=0,1,\cdots,\gamma N$、

$\mathrm{div}$ と $\mathrm{mod}$ は第一引数を第二引数で割ったときの商と余りを計算する関数をそれぞれ表して

いる。この場合、$y$ の信号長は $\gamma N$ となる。また、$x$ を倍率 $\gamma$ でダウンサンプリングした結果を

$z$ とすると、$z$ の信号長は $N/\gamma$ となり以下の式で定義されている。

\begin{equation}
z_n = x_{n\times\gamma}
\tag{5}
\end{equation}

アップサンプリングとダウンサンプリング後の信号 $y$ と $z$ を DFT するとそれぞれ $X(\gamma\omega)$

と $X(\omega/\gamma)$ となる。アップサンプリングの場合は、$Y(1)$ の周波数特性が $X(\gamma)$ と等価となるので、

元の信号の周波数特性 $X(\omega)$ が周波数軸上において倍率 $\gamma$ で縮められた周波数特性に変換された

とみなせる。一方ダウンサンプリングの場合は、周波数軸上において倍率 $\gamma$ で引き延ばされたとみ

なせる。つまり、時間軸上で引き延ばされる処理は周波数軸上で縮められており、逆もまた同様で

ある。また、2.2 項で示したことと同様に、二次元信号の場合は、縦横それぞれの軸に対して上記

処理が可能であり、周波数解析もそれらの軸に沿って収縮と拡張がされる。

\subsection{離散コサイン変換}

離散コサイン変換（DCT: Discrete cosine transform）は、DFT の一種であり、実数値での高速
処理が実現できるため、画像符号化方式の JPEG 等で広く用いられている。2.2 で示したとおり、
周波数解析には DFT を用いるが複素数値の処理となるため、計算量が多く時間がかかるのが欠点
である。DCT は、対象となる信号を対称に拡張しそれに DFT を適用して、対称性を考慮して展
開すると虚部が打ち消し合うことから実数値演算として導出される。DCT は、拡張の仕方により
離散サイン変換も含めて四タイプが存在しているが、広く用いられているのはタイプ 2 のものであ
る。タイプ 2 の DCT の変換行列 $D$ は以下の式で定義されている。

\begin{equation}
d_{ij} =
\sqrt{\frac{2}{M}}
c_i c_j
\cos
\left(
\frac{i(j+1/2)\pi}{M}
\right)
\tag{6}
\end{equation}

ここで、$d_{ij}$ は $D$ の $(i,j)$ 要素（$i,j=0,1,\cdots,M-1$）を表し、$M$ は $D$ のサイズ、
$c_k$ は係数をそれぞれ表している。$c_k$ は、$D$ が正規直交行列となるために、以下のように定義
されている。

\[
c_k=
\begin{cases}
1/\sqrt{2} & \mathrm{if}\ k=0 \\
1 & \mathrm{otherwise}
\end{cases}
\]

よって、$D^T D = I$ である。ここで、上付きの $T$ は転置を表し、$I$ は単位行列をそれぞれ表して
いる。サンプル数 $M$ の一次元信号 $x$ を DCT するには、$y=Dx$ とする。ここで、$x$ は列ベクトル
とする。$y$ の各要素値は、$x$ のそれぞれ異なるある帯域における周波数成分の値を表している。つ
まり、$D$ の行ベクトル $d_i$ は等長かつ異なる通過域を有するフィルタと考えられる。例えば
$M=4$ の場合、各 $d_i$ の通過域は $[0,\pi]$ を四分割したものとなる。よって $M$ は、チャンネル数
と呼ばれており、周波数解析における分解能を示している。

画像 $X$ に DCT を適用することは以下の式で表せる。

\begin{equation}
Y=\Lambda_H X \Lambda_W^T
\tag{7}
\end{equation}

ここで、$Y$ は変換後画像を表しており、$H$ と $W$ は $X$ の縦と横の長さを、$\Lambda_k$ は長さ $k$
であり $D$ が要素であるブロック対角行列をそれぞれ表している。ブロック対角行列とは、正方行
列であり、その対角要素に同じ行列が並んでいるものである。そのため、DCT を適用する場合は
画像の長さを DCT の $M$ の倍数となるように拡張する必要がある。式(7)において、$\Lambda_H$ は
$X$ の縦方向における周波数変換を担っており、$\Lambda_W$ は横方向である。$Y$ は、$M\times M$
のブロック毎に左上から右下にかけて低周波数成分から高周波数成分の値が並んでおり、それが縦
と横それぞれ $H/M$ 個と $W/H$ 個並んでいる。その表示状態をブロックモードという。$Y$ を各
周波数帯域成分毎に $H/M \times W/H$ の行列にまとめて、それら行列を元の順番に沿って配置し
直したものをサブバンドモードという。

式(7)の逆変換は、$D^T D = I$ より、以下の式で表せる。

\begin{equation}
\Lambda_H^T Y \Lambda_W
=
\Lambda_H^T
(
\Lambda_H X \Lambda_W^T
)
\Lambda_W
=
(
\Lambda_H^T \Lambda_H
)
X
(
\Lambda_W^T \Lambda_W
)
=
X
\tag{8}
\end{equation}

\subsection{低域通過フィルタ}

フィルタリングにおいて直流成分を含めた低周波数帯域の成分のみを出力するものを低域通過
フィルタと呼ぶ。画像に適用される二次元フィルタは縦と横の方向それぞれに対して周波数特性が
存在し、二次元の低域通過フィルタでは縦と横方向それぞれの低周波数成分の値が出力される。ガ
ウシアンフィルタ $G$ は、最も有名な低域通過フィルタであり、以下の式で定義されている。

\begin{equation}
g_{ij}
=
\frac{1}{2\pi\sigma^2}
\exp
\left(
-\frac{i^2+j^2}{2\sigma^2}
\right)
\tag{9}
\end{equation}

ここで、$g_{ij}$ は $G$ の $(i,j)$ 要素（$i\in[-k_i,k_i]\in\mathbb{Z},j\in[-k_j,k_j]\in\mathbb{Z}$）
を表し、$\sigma>0$ はフィルタ係数をコントロールするパラメータであり、式(1)と同様に $k_i$
と $k_j$ はフィルタ長をそれぞれ表している。原点は直流成分に対する特性を表し、端に行くほど
各方向における高周波数成分に対する特性を表している。


\subsection{高域通過フィルタ}

高域通過フィルタは、低域通過フィルタと反対に、高周波数帯域の成分のみを出力するものであ
る。画像に適用される場合、縦と横方向のどちらに対しての高周波数成分を出力させるのか考慮す
る必要がある。ソーベルフィルタ $S$ は、一次微分を行うフィルタであり、画像に適用する場合には
以下のとおり縦と横方向の高域通過フィルタになるものが一般的である。

\[
\begin{bmatrix}
-1 & -2 & -1 \\
0 & 0 & 0 \\
1 & 2 & 1
\end{bmatrix}
\qquad
\begin{bmatrix}
-1 & 0 & 1 \\
-2 & 0 & 2 \\
-1 & 0 & 1
\end{bmatrix}
\]

ここで、左のフィルタ行列が縦方向のものであり、右のフィルタ行列が横方向である。

またラプラシアンフィルタ $L$ は、同様に二次微分を行うフィルタであり、縦と横方向ともに高周
波数である成分を出力する。画像に適用する場合、一般的には以下にしめす $3\times3$ のフィルタ行列
サイズを有する四近傍型と八近傍型が用いられる。

\[
\begin{bmatrix}
0 & 1 & 0 \\
1 & -4 & 1 \\
0 & 1 & 0
\end{bmatrix}
\qquad
\begin{bmatrix}
1 & 1 & 1 \\
1 & -8 & 1 \\
1 & 1 & 1
\end{bmatrix}
\]

ここで、左と右はそれぞれ四近傍型と八近傍型であり、基本的な性質は同じであるものの周波数特
性が異なる。

\subsection{画像のノイズ}

画像には、様々な要因によってノイズが混入することがあり、その除去には平滑化処理が古典的
な手法として有効である。ノイズの要因は、撮影対象からの光情報を電子データとして保存するま
での過程における微小な誤差であり、最も影響が大きいものは受光素子上での熱雑音である。ノイ
ズがある撮影画像 $Y$ は一般的に以下の式で定義されている。

\begin{equation}
Y = X + N
\tag{10}
\end{equation}

ここで、$X$ と $N$ はノイズがない原画像とノイズ成分をそれぞれ表している。$N$ の要素値は、各要
素間において独立であり、ある確率密度関数に従う乱数としてモデル化されるのが一般的である。
モデルは、様々なものが議論されているが、最も汎化性が高いものとしてガウス分布がしばしば用
いられる。ガウス分布 $N(\mu,\sigma^2)$ の確率密度関数は以下の式で定義されている。

\begin{equation}
f(x)
=
\frac{1}{\sqrt{2\pi\sigma^2}}
\exp
\left(
-\frac{(x-\mu)^2}{2\sigma^2}
\right)
\tag{11}
\end{equation}

ここで、$\mu$ と $\sigma^2$ はそれぞれ分布の平均値と分散値を表している。画像のノイズの場合は
$\mu=0$ とするのが一般的である。そのため、ノイズ除去として 2.5 節で示した低域通過フィルタ
による平滑化が有効となる。平滑化により、$N$ は平均値つまり 0 に漸近していくので低減する。
ただし、平滑化後の画像は原画像に平滑化が適用されたものとなるので、全体的にぼやけた画像と
なる欠点がある。

\subsection{画像の先鋭化}

先鋭化とは、画像の濃淡を残しながらエッジ部分を強調する技術であり、古典的な手法としてア
ンシャープマスキングが提案されている。アンシャープマスキングでは、平滑化を用いて入力画像
からエッジ成分を算出し、エッジ成分を増幅して入力画像に加算することでエッジ強調を実現して
いる。アンシャープマスキングは以下の式で定義されている。

\begin{equation}
Y = X + p(X-g(X))
\tag{12}
\end{equation}

ここで、$X$ と $Y$ は入力画像と強調後画像を表しており、$p$ は強調の強さを調整するパラメータ、
$g$ は平滑化を行う関数をそれぞれ表している。アンシャープマスキングでは、$X$ から平滑化を適用
した $g(X)$ を減算してエッジ成分を算出しており、それを $p$ 倍して入力画像に加算することで、
$X$ の全体的なバランスを維持しつつエッジ部分の強調を実現している。
\section{方法}

\subsection{一次元信号の周波数解析}

周波数 $\alpha =$ をサイン関数に適用し一次元信号（横ベクトル）$x$ を作成する。次に、作成した
一次元信号を縦方向に引き延ばして正方の画像 $X$ に拡張する。最後に、DFT を用いて $x$ の周波
数振幅特性を計算しグラフ $A$ に表示する。

また、2.3 節を基にして、$X$ を倍率 $\gamma =$ でアップサンプリングを行い $X_{up}$ を作成する。
また、$\gamma$ でダウンサンプリングを行い $X_{down}$ を作成する。上記と同様に、それらから一次元
信号 $x_{up}$ と $x_{down}$ を抜き出し、DFT を用いて周波数振幅特性を計算しグラフ $A_{up}$ と
$A_{down}$ に表示する。

また、上記処理を $\alpha =$ と $\gamma =$ 、$\alpha =$ と $\gamma =$ において行う。

\subsection{二次元信号の周波数解析}

画像 を読み込み $X$ とする。次に、式(6)を基にして、$M=4$ の DCT 変換行列 $D$ を作成す
る。最後に、式(7)と $D$ を用いて $X$ を変換し、変換後の画像 $Y$ を算出する。ここで、$Y$ はブロッ
クモードのままであるのでサブバンドモードに変換する。

また、上記処理を画像 と画像 において行う。

\subsection{画像の逆変換}

[この節は達成した者のみが任意で記述する。つまり省略しても構わない。]

3.2 節の変換後の画像 $Y$ に対して、式(8)に沿って逆変換を行い、逆変換後の画像 $Z$ を算出す
る。また、以下の式で入力画像 $X$ と $Z$ の誤差 $L$ を計算する。

\begin{equation}
L(X,Z)
=
\sum_{i=0}^{H}
\sum_{j=0}^{W}
|x_{ij}-z_{ij}|
\tag{13}
\end{equation}

ここで、$H$ と $W$ はそれぞれ $X$ の縦と横の長さを表し、$x_{ij}$ と $z_{ij}$ はそれぞれ $X$ と
$Z$ の $(i,j)$ 要素を表している。

\subsection{ガウシアンフィルタによる画像のぼかし}

式(9)を基にして、フィルタ長 $k =$ と標準偏差 $\sigma =$ を用いて一次元のガウシアンフィルタ
$g$ を作成する。また、DFT を用いて $g$ の周波数振幅特性を計算しグラフ $G$ に表示する。次に、
$k$ と $\sigma$ を用いて二次元のガウシアンフィルタ $G$ を作成する。最後に、画像 を読み込み $X$
とし、$X$ に $G$ を畳み込んでフィルタリング後の出力画像 $Y$ を算出する。

また、上記処理を以下の組み合わせで行う

\begin{itemize}
\item $k =$ 、$\sigma =$ 、画像 。
\item $k =$ 、$\sigma =$ 、画像 。
\item $k =$ 、$\sigma =$ 、画像 。
\end{itemize}

\subsection{画像のノイズとその除去}

画像 を読み込んで $X$ とし、画像の縦と横の画素数を確認する。次に、式(11)（$\mu=0$）を基
にして、ノイズの標準偏差 $\sigma_n =$ を用いてノイズ成分 $N$ を算出し、ノイズあり画像
$Y=X+N$ を算出する。最後に、3.4 節と同様の処理を行い、平滑化後の画像 $Z$ を算出する。

また、上記処理を以下の組み合わせで行う

\begin{itemize}
\item $k =$ 、$\sigma =$ 、$\sigma_n =$ 、画像 。
\item $k =$ 、$\sigma =$ 、$\sigma_n =$ 、画像 。
\item $k =$ 、$\sigma =$ 、$\sigma_n =$ 、画像 。
\item $k =$ 、$\sigma =$ 、$\sigma_n =$ 、画像 。
\end{itemize}

\subsection{ソーベルフィルタによる画像のエッジ検出}

まず、2.6 節を基にして、縦と横方向のソーベルフィルタ行列 $S_v$ と $S_h$ をそれぞれ作成する。
ここで、DFT を用いて $S_h$ の第一列ベクトルの周波数振幅特性を計算しグラフ $S$ に表示する。
最後に、画像 を読み込んで $X$ とし、$X$ に $S_v$ と $S_h$ をそれぞれ畳み込み、フィルタリング後
の出力画像 $Y_v$ と $Y_h$ を算出する。

また、上記処理を画像 と画像 において行う。

\subsection{ラプラシアンフィルタによる画像のエッジ検出}

2.6 節を基にして、八近傍型のラプラシアンフィルタ行列 $L$ を作成する。次に、任意の画像 $X$
を読み込んで $L$ を畳み込み、フィルタリング後の出力画像 $Y$ を算出する。それら $X$ と $Y$ を png
画像として保存する。

また、ガウシアンフィルタを用いて同様に実現する。まず、3.4 節と同様にガウシアンフィルタ
$G$ を作成する。次に、画像 を読み込んで $X$ とし、$X$ に $G$ を畳み込み、フィルタリング後の出
力画像 $Y$ を算出する。最後に、$Z=X-Y$ を計算する。その $Z$ を png 画像として保存する。

また、上記二つの処理を以下の組み合わせで行う

\begin{itemize}
\item $k =$ 、$\sigma =$ 、画像 。
\item $k =$ 、$\sigma =$ 、画像 。
\item $k =$ 、$\sigma =$ 、画像 。
\end{itemize}

\subsection{アンシャープマスキングによる画像のエッジ強調}

[この節は達成した者のみが任意で記述する。つまり省略しても構わない。]

2.8 節を基にしてアンシャープマスキングによるエッジ強調を実現する。まず、3.4 節と同様に
ガウシアンフィルタ $G$ を作成する。次に、画像 を読み込んで $X$ とし、$X$ に $G$ を畳み込み、
フィルタリング後の出力画像 $\hat{Y}$ を算出する。次に、$Y=X-\hat{Y}$ を計算してエッジ成分
$Y$ を算出する。最後に、エッジ強調のパラメータ $p =$ を用いて強調後の出力画像
$Z=X+pY$ を算出する。それら $X$ と $Y$、$Z$ を png 画像として保存する。

また、上記処理を以下の組み合わせで行う

\begin{itemize}
\item $k =$ 、$\sigma =$ 、$p =$ 、画像 。
\item $k =$ 、$\sigma =$ 、$p =$ 、画像 。
\item $k =$ 、$\sigma =$ 、$p =$ 、画像 。
\end{itemize}








\section{結果と考察}

\subsection{一次元信号の周波数解析}

ここに、3.1 節の結果と考察を示す。図 1 に、$\alpha =$ における画像 $X$ と周波数振幅特性のグラ
フ $A$ をそれぞれ示す。また図 2 に、$\gamma =$ においてアップサンプリングとダウンサンプリング
した画像 $X_{up}$ と $X_{down}$、それらの周波数振幅特性のグラフ $A_{up}$ と $A_{down}$ をそれぞれ
示す。[これらの図に対する考察を述べる。]

[同様に、他の設定における結果の表示とそれに対する考察を述べる。]

% \begin{figure}[htbp]
% \centering

% \begin{minipage}{0.45\linewidth}
% \centering
% \includegraphics[width=\linewidth]{fig1a.png}
% \\
% (a) X
% \end{minipage}
% \hfill
% \begin{minipage}{0.45\linewidth}
% \centering
% \includegraphics[width=\linewidth]{fig1b.png}
% \\
% (b) A
% \end{minipage}

% \caption{サイン関数による画像と周波数振幅特性}
% \end{figure}

% \begin{figure}[htbp]
% \centering

% \begin{minipage}{0.22\linewidth}
% \centering
% \includegraphics[width=\linewidth]{fig2a.png}
% \\
% (a) $X_{up}$
% \end{minipage}
% \hfill
% \begin{minipage}{0.22\linewidth}
% \centering
% \includegraphics[width=\linewidth]{fig2b.png}
% \\
% (b) $A_{up}$
% \end{minipage}
% \hfill
% \begin{minipage}{0.22\linewidth}
% \centering
% \includegraphics[width=\linewidth]{fig2c.png}
% \\
% (c) $X_{down}$
% \end{minipage}
% \hfill
% \begin{minipage}{0.22\linewidth}
% \centering
% \includegraphics[width=\linewidth]{fig2d.png}
% \\
% (d) $A_{down}$
% \end{minipage}

% \caption{アップサンプリングとダウンサンプリング後の画像と周波数振幅特性}
% \end{figure}

\subsection{二次元信号の周波数解析}

ここに、3.2 節の結果と考察を示す。図 3 に、入力画像 $X$ とその DCT を適用した後の出力画像
$Y$ をそれぞれ示す。[これらの図に対する考察を述べる。]

[同様に、他の設定における結果の表示とそれに対する考察を述べる。]

% \begin{figure}[htbp]
% \centering

% \begin{minipage}{0.45\linewidth}
% \centering
% \includegraphics[width=\linewidth]{fig3a.png}
% \\
% (a) X
% \end{minipage}
% \hfill
% \begin{minipage}{0.45\linewidth}
% \centering
% \includegraphics[width=\linewidth]{fig3b.png}
% \\
% (b) Y
% \end{minipage}

% \caption{入力画像と DCT 後の出力画像}
% \end{figure}

\subsection{画像の逆変換}

[この節は 3.3 節を記述した者は必須。]

ここに、3.3 節の結果と考察を示す。図 4 に、入力画像 $X$ とそれを DCT した後に逆変換を行っ
た出力画像 $Z$ をそれぞれ示す。また、これら $X$ と $Z$ の絶対値総和誤差 $L$ は であった。
[これらの図に対する考察を述べる。]

% [同様に、他の設定における結果の表示とそれに対する考察を述べる。]

% \begin{figure}[htbp]
% \centering

% \begin{minipage}{0.45\linewidth}
% \centering
% \includegraphics[width=\linewidth]{fig4a.png}
% \\
% (a) X
% \end{minipage}
% \hfill
% \begin{minipage}{0.45\linewidth}
% \centering
% \includegraphics[width=\linewidth]{fig4b.png}
% \\
% (b) Z
% \end{minipage}

% \caption{入力画像と DCT の逆変換後の出力画像}
% \end{figure}

\subsection{ガウシアンフィルタによる画像のぼかし}

ここに、3.4 節の結果と考察を示す。図 5 に、$k =$ と $\sigma =$ におけるガウシアンフィルタ
の周波数振幅特性のグラフ $G$ と入力画像 $X$、出力画像 $Y$ をそれぞれ示す。
[これらの図に対する考察を述べる。]

[同様に、他の設定における結果の表示とそれに対する考察を述べる。]

% \begin{figure}[htbp]
% \centering

% \begin{minipage}{0.3\linewidth}
% \centering
% \includegraphics[width=\linewidth]{fig5a.png}
% \\
% (a) G
% \end{minipage}
% \hfill
% \begin{minipage}{0.3\linewidth}
% \centering
% \includegraphics[width=\linewidth]{fig5b.png}
% \\
% (b) X
% \end{minipage}
% \hfill
% \begin{minipage}{0.3\linewidth}
% \centering
% \includegraphics[width=\linewidth]{fig5c.png}
% \\
% (c) Y
% \end{minipage}

% \caption{ガウシアンフィルタの周波数振幅特性と画像に適用した場合の入力と出力画像}
% \end{figure}





\end{document}